\documentclass[11pt,a4paper]{article}
\usepackage[T1]{fontenc}
\usepackage{mathptmx}
\usepackage[margin=25mm]{geometry}
\usepackage{booktabs}
\usepackage{array}
\usepackage{amsmath}
\usepackage{microtype}
\usepackage{titlesec}
\usepackage[hidelinks]{hyperref}
\usepackage{fancyhdr}
\titleformat{\section}{\large\bfseries}{\thesection}{0.6em}{}
\titleformat{\subsection}{\bfseries}{\thesubsection}{0.6em}{}
\setlength{\parskip}{0.5em}
\setlength{\parindent}{0pt}
\pagestyle{fancy}\fancyhf{}
\renewcommand{\headrulewidth}{0.4pt}
\newcommand{\inr}[1]{INR~#1}

\fancyhead[L]{RescueSwarm --- Track C: Autonomy and Ground Control}
\fancyhead[R]{\thepage}
\begin{document}

\begin{center}
{\LARGE\bfseries Track C --- Autonomy and Ground Control}\\[3mm]
{\large RescueSwarm: work package and funding request}\\[2mm]
NIDAR 2026--27 $\cdot$ Track 1 $\cdot$ Mission 1\\[1mm]
\textbf{Ask: \inr{2,324}} $\cdot$ 2 students $\cdot$ Coverage planner, task allocation, state machine, ground station, SITL
\end{center}
\vspace{2mm}\hrule\vspace{4mm}


This track asks for the least money in the programme and carries a large share
of the scored functionality. That is not an accident of accounting: the
deliverables are software, the workstations are team-supplied, and the
verification environment is free. It is worth stating explicitly, because a
funding split that shows a near-zero line invites the assumption that the track
is small. It is not --- it is inexpensive.

\section{Scope and ownership}

This track owns \textbf{coverage planner, task allocation, state machine, ground station, sitl}. It is one of four delivery tracks defined in
\texttt{docs/development-plan.md}~\S1.3; the integrated system argument lives in
the master proposal, \texttt{docs/proposal/rescueswarm-proposal.pdf}, which this
document does not restate.

\section{Coverage planning over arbitrary search regions}

The rulebook does not promise a rectangular search area, so the planner does not
assume one. The implementation decomposes a possibly non-convex boundary into
monotone cells, fills each by scanline, clips every transect to the true
boundary, and routes between cells along a visibility graph shortest path rather
than a straight line that might leave the region.

Verified across 13 boundary shapes including non-convex notches and slots:
\textbf{0.00\,m of path outside the boundary}, workload imbalance across three
aircraft \textbf{$\leq$0.02\,\%}, and 157 tests passing. Three defects were found
and fixed in reaching that state --- a scanline span that bridged concave
notches, a missing cell decomposition, and a straight inter-cell hop that cut
corners across excluded ground.

\section{Autonomy, and how it is evidenced}

The competition requires autonomous operation, and asserting it is not the same
as showing it. The verification harness parses the flight script's own syntax
tree and proves that \textbf{no MAVLink transmission occurs after the mode is
set to AUTO}. That is a structural argument rather than an observational one:
it does not depend on having watched a particular flight behave.

The same harness checks inter-aircraft separation, geofence containment,
failsafe configuration, energy consumption against the model, and recorder clock
integrity, and exits non-zero on any failure. It is re-runnable, which means the
autonomy claim is regenerated rather than remembered.

A separation breach of 3.10\,m was found this way and corrected to 5.34\,m by
staggering recovery loiter times. A 453.9\,s recorder clock discontinuity was
found the same way, and independently confirmed against consumed charge.

\section{The ground station is deliberately unable to fly the aircraft}

Requirement SYS-20 holds that the ground station must be \emph{architecturally}
incapable of originating a retask, waypoint change or arming command --- not
merely configured not to. This is a design property, not a policy, and it is
what makes the autonomy claim defensible under scrutiny: an operator cannot
accidentally invalidate it.

\section{Why the ask is small}

The two workstations are team-supplied and the simulation environment is
open-source. What remains is a fabricated sun hood and an observer monitor for
field use. This track's real cost is student time across the full programme
duration, which the funding request does not price and which the schedule must
protect.

\section{Funding request}

Every figure below is generated from
\texttt{docs/proposal/figures/track\_budget.py}, which partitions the master
competition budget. The four track asks plus the shared pool reconcile to the
master figure exactly; that reconciliation is asserted in code and fails the
build if it ever stops holding. \textbf{K} marks a line held at professional
grade on purpose.


\begin{center}
\begin{tabular}{@{}p{88mm}cr@{}}
\toprule
\textbf{Item} & & \textbf{\inr{}} \\
\midrule
Sun hood + observer monitor &  & 1,500 \\
\midrule
Subtotal & & 1,500 \\
Customs duty and freight & & 213 \\
GST & & 308 \\
Contingency at 15\% & & 303 \\
\midrule
\textbf{TRACK C ASK} & & \textbf{2,324} \\
\bottomrule
\end{tabular}
\end{center}


\textbf{Deferred or institutionally held, and therefore not requested:}
Backup GCS laptop, GCS laptop. These remain capabilities the track requires; they are simply not being
bought. Where a deferral carries a technical consequence it is stated in the
risk table below rather than left implicit.

\section{Deliverables by phase}

\begin{center}
\begin{tabular}{@{}llp{88mm}@{}}
\toprule
\textbf{Phase} & \textbf{Dates} & \textbf{This track delivers} \\
\midrule
\textbf{P1} & 24 Aug -- 13 Sep & SITL flying one aircraft. Planner validated against arbitrary boundaries. \\
\textbf{P2} & 14 Sep -- 4 Oct & Full mission in SITL. Verification harness passing. \\
\textbf{P3} & 5 Oct -- 18 Oct & Simulation re-baselined against the built aircraft. \\
\textbf{P5} & 9 Nov -- 22 Nov & Autonomous full mission, no operator input after AUTO. \\
\bottomrule
\end{tabular}
\end{center}

\section{Risks owned by this track}

\begin{center}
\begin{tabular}{@{}p{50mm}lp{70mm}@{}}
\toprule
\textbf{Risk} & \textbf{Sev.} & \textbf{Mitigation} \\
\midrule
Pad containment has zero margin & High & Recorded in the master proposal as a measured result. Recovery geometry must be re-verified whenever aircraft mass or descent rate changes. \\
Survey altitude is unresolved & Medium & The design point says 60\,m; every flown simulation uses 40\,m. Coverage timing, transect count and detection all move with it. Resolution is shared with Track D and needs a recall measurement to settle. \\
Simulation diverges from the built aircraft & Medium & SITL evidence is only evidence if the modelled aircraft matches the built one. Re-baseline against Track A's measured mass and thrust after P2. \\
\bottomrule
\end{tabular}
\end{center}

\section{Interfaces to other tracks}

\begin{center}
\begin{tabular}{@{}lp{110mm}@{}}
\toprule
\textbf{With} & \textbf{Dependency} \\
\midrule
Track A & Consumes measured mass, thrust and endurance. Simulation fidelity depends on receiving them after the P2 thrust-stand run. \\
Track B & Runs on the companion computer and commands through its MAVLink routing. Depends on the link carrying control and telemetry together. \\
Track D & Consumes detections and supplies the aircraft state --- position, attitude, time --- that each geotag is computed from. \\
\bottomrule
\end{tabular}
\end{center}

Interfaces are where student programmes fail, because each side assumes the
other owns the gap. Every row above is a two-way commitment and should be
recorded as an interface control document at P1.


\vfill
\hrule\vspace{2mm}
{\small Generated by \texttt{tools/proposal/build\_track\_proposals.py} from
\texttt{docs/proposal/figures/track\_budget.py}. Financial figures are not
maintained in this document and must not be edited here: change the master
budget and regenerate. Technical claims cite the model or document that
produced them.}
\end{document}
